\section{\uppercase{Introduction}}
\label{sec:introduction}
\begin{figure*}[t]
\centering
\resizebox{\textwidth}{!}{%
\begin{tikzpicture}[
  font=\footnotesize,
  stage/.style={draw=black!40, rounded corners=2pt, fill=#1, align=left, inner sep=4pt, minimum height=3.1cm, anchor=north west},
  exper/.style={draw=black!40, rounded corners=2pt, fill=#1, align=left, inner sep=4pt, text width=5.35cm, minimum height=2.55cm, anchor=north west},
  arr/.style={-{Stealth[length=2mm]}, thick, black!55},
]
\coordinate (bus) at (0,-3.5);
\coordinate (row2) at (0,-3.9);

\node[stage=blue!4, text width=3.6cm] (data) at (0,0)
  {\textbf{PaHaW recordings}\\[2pt]
   75 participants, 8 tasks,\\597 recordings\\[3pt]
   per sample: position, time, pen contact, pressure, altitude and azimuth};

\node[stage=orange!7, text width=4.0cm, right=5mm of data.north east, anchor=north west] (render)
  {\textbf{Pseudodynamic minimap}\\[2pt]
   \textbullet\ trace drawn only where the pen touches the tablet\\
   \textbullet\ each signal normalized within its own recording\\
   \textbullet\ azimuth centered on its circular mean\\
   \textbullet\ up to three signals mapped to RGB};

\node[stage=orange!7, text width=3.9cm, right=5mm of render.north east, anchor=north west] (enc)
  {\textbf{Nine encodings} (Table~\ref{tab:encoding})\\[3pt]
   \begin{tabular}{@{}c@{\hspace{8pt}}c@{}}
   \includegraphics[height=1.2cm]{method_static.png} & \includegraphics[height=1.2cm]{method_saz.png}\\
   \scriptsize 1: static trace & \scriptsize 8: $SAZ$
   \end{tabular}\\[2pt]
   2--5: one signal; 6--9: RGB triplets};

\node[stage=blue!4, text width=3.7cm, right=5mm of enc.north east, anchor=north west] (eval)
  {\textbf{Evaluation}\\[2pt]
   $5\times5$ participant-separated cross-validation with 3 inner folds\\[3pt]
   recording-level macro F1\\[1pt]
   participant-level macro F1 (majority vote over tasks)};

\node[exper=green!6] (e1) at (data.north west |- row2)
  {\textbf{Experiment 1} \hfill Sec.~\ref{sec:descriptor}--\ref{sec:classifiers}\\[2pt]
   all eight tasks\\[2pt]
   \textbullet\ each encoding described by LPQ texture and classified by SVM\\
   \textbullet\ leading encoding ($SAZ$): SVM compared with logistic regression, random forest and a CNN trained from scratch};

\node[exper=green!6, right=5mm of e1.north east, anchor=north west] (e2)
  {\textbf{Experiment 2} \hfill Sec.~\ref{sec:nested-followup}--\ref{sec:controls}\\[2pt]
   all eight tasks\\[2pt]
   \textbullet\ encoding, SVM and fusion rule chosen with training participants only\\
   \textbullet\ spatial controls: signals circularly shifted or permuted within each stroke};

\node[exper=green!6, right=5mm of e2.north east, anchor=north west] (e3)
  {\textbf{Experiment 3} \hfill Sec.~\ref{sec:competitive}\\[2pt]
   tasks T2--T4 (all eight as secondary)\\[2pt]
   \textbullet\ static trace + LPQ + SVM\\
   \textbullet\ minimaps + LPQ + SVM\\
   \textbullet\ minimaps + ResNet18 features\\
   \textbullet\ kinematic/pressure attributes + SVM};

\draw[arr] ([yshift=-1.55cm]data.north east) -- ([yshift=-1.55cm]render.north west);
\draw[arr] ([yshift=-1.55cm]render.north east) -- ([yshift=-1.55cm]enc.north west);
\draw[arr] ([yshift=-1.55cm]enc.north east) -- ([yshift=-1.55cm]eval.north west);

\draw[thick, black!55] (eval.south) -- (eval.south |- bus);
\draw[thick, black!55] (e1.north |- bus) -- (eval.south |- bus);
\foreach \e in {e1,e2,e3} \draw[arr] (\e.north |- bus) -- (\e.north);

\draw[arr, dashed, black!45] (data.north) -- ++(0,0.45)
  -| ([xshift=5mm,yshift=-2.05cm]e3.north east)
  node[pos=0.25, above, font=\scriptsize, text=black!70] {pen signals used directly, without an image}
  -- ([yshift=-2.05cm]e3.north east);
\end{tikzpicture}}
\caption{Overview of the method. Each PaHaW recording is rendered as a pseudodynamic minimap in one of nine encodings, and every experiment uses the same participant-separated cross-validation. Experiment~1 compares encodings and classifiers, Experiment~2 chooses the encoding with training participants only and tests whether the position of the signals along the stroke matters, and Experiment~3 compares minimaps with the static trace and with kinematic/pressure attributes computed directly from the pen signals (dashed arrow). Thumbnails show the spiral of one healthy participant as rendered.}
\label{fig:method}
\end{figure*}

\noindent
A digitizing tablet records the pen trajectory together with its speed, pressure and inclination. Encoding these measurements as colour adds information about how a trace was produced while keeping its recognizable shape, so that standard image descriptors and classifiers can use the dynamics directly.

The application is the distinction between Parkinson's disease (PD) and healthy participants (H) in PaHaW, a database of online handwriting and drawing tasks \cite{Drotar2016}. Handwriting exposes aspects of fine motor behaviour, and the relation between pen motion and diagnosis depends on the task, device and participant characteristics \cite{Impedovo2019}.

We propose \emph{pseudodynamic minimaps}, a new image representation of a handwriting task execution. Planar position supplies the image coordinates and colour carries the measured pen attributes. The term minimap follows work that renders source code compactly as a visual signature \cite{Gebara2025}.

Handwriting classification has often used explicitly designed kinematic and pressure features, as in the study that introduced PaHaW \cite{Drotar2016}. Learned features provide another route, either from pen dynamics \cite{Pereira2016} or from visual attributes of reconstructed writing \cite{Moetesum2019}. Dynamic information has also appeared in earlier handwriting images: Diaz et al. encode movement through acquired points and pen-up trajectories for PD detection \cite{Diaz2019}, and Cilia et al. build synthetic RGB images from handwriting dynamics for Alzheimer's disease detection \cite{Cilia2021SyntheticImages}.

Rendering a temporal signal as an image is an established route to reusing visual descriptors and convolutional models. Recurrence plots expose the recurrence structure of a dynamical system as a two-dimensional pattern \cite{Eckmann1987}, Gramian angular fields and Markov transition fields encode a univariate series as a matrix \cite{Wang2015}, and convolutional networks trained on such renderings are competitive on time-series benchmarks \cite{Hatami2018}. The same reduction appears outside the time domain, where source code is rendered as an image for classification \cite{Shi2022}.

Pseudodynamic minimaps differ from these transforms in three respects. The image coordinates are the physical positions the pen visited rather than a time axis, so each signal value stays anchored to the stroke that produced it. Each signal is normalized within its own recording by percentiles, so colour describes the writer's relative variation. Azimuth is centred on its circular mean, so that the arbitrary zero-degree boundary does not separate similar orientations.

This initial study evaluates the representation in three experiments that share participant-separated partitions (Figure~\ref{fig:method}). Experiment~1 adds speed, pressure, altitude and azimuth to the static trace, one at a time and as RGB triplets, and compares classifiers on the leading encoding. Experiment~2 selects the encoding with training participants only and displaces the signals along each stroke to test whether their position matters. Experiment~3 compares minimaps, described by LPQ texture or by pretrained ResNet18 features, with the static trace and with kinematic/pressure attributes on the repeated \textit{l}, \textit{le} and \textit{les} tasks examined in earlier work on PaHaW \cite{Casademunt2023}.

The results indicate that the classifiers make use of the encoded dynamics. Training-only selection preferred a dynamic encoding over the static trace in almost every training set, displacing the signals along the stroke lowered performance, and on the three repeated-pattern tasks the minimap pipeline gave the best observed participant-level decisions among the four pipelines. A cohort of 75 participants leaves the size of these gains uncertain; the study establishes how pen dynamics can be carried into an image and where they help, rather than delivering a clinical tool.
